\chapter{Introduction}
\label{ch:introduction}

Cities are nexuses of ideas, culture, and opportunity, housing and sustaining more than half of the global population. However, urban expansion represents one of the most irreversible and environmentally damaging forms of land conversion in human history, depleting freshwater resources, exacerbating flooding, diminishing biodiversity, and creating urban heat islands with unprecedented temperature extremes~\cite{gao2020Mapping,zhang2018Urbanization}. The identification, repurposing, and intentional use of \acrfull{uvl} is essential to mitigate further land consumption while satisfying the needs of expanding urban populations. This is particularly important in the American Northeast megalopolis, where land constraints, flooding risk, intensifying housing crises, and escalating costs of living are culminating in population decline.

\section{Land Consumption in America}
\label{sec:land_consumption}

While the UN and several EU member states have established targets to reduce or achieve net-zero land consumption by 2050, the United States has no comparable policy and is widely associated with urban sprawl~\cite{xu2022Automatica,barrington-leigh2015century}. Despite the environmental impacts of urban expansion, development in the U.S.\ continues to perpetuate low-density, car-centric land-use patterns introduced by postwar policy initiatives favoring detached single-family housing. These trends hollowed out dense urban cores to create suburban networks that are difficult to retrofit toward density due to entrenched low-connectivity road networks~\cite{barrington-leigh2015century}. To reduce American urban land consumption and curb its environmental effects, the most feasible path is to densify and intelligently utilize existing urban cores through infill development. Furthermore, strategically repurposing \acrshort{uvl} can support both increased housing supply and environmental resilience. However, identifying \acrshort{uvl} is a costly and manual process, particularly in jurisdictions with limited digitized cadastral data and planning resources. An open-source pipeline to identify \acrshort{uvl} would provide a vital tool to cities facing both housing and environmental challenges.

\section{UVL for the Housing Crisis}
\label{sec:housing_crisis}

The U.S.\ is facing a shortage of 4.7 million homes, stemming from a decade of conservative development during the Great Recession and contributing to the current affordability crisis~\cite{myers2025Misalignment}. Increasing the supply of homes is essential to addressing the ongoing housing crisis. Unfortunately, in the U.S., buyer-friendly, elastic housing markets like those in the Sun Belt are inextricably linked to landlocked regions with the ability to sprawl~\cite{saiz2010geographic,jysma2017Geographic,vogler2021Trends}. The identification and redevelopment of \acrshort{uvl} to housing would increase housing supply elasticity without contributing to sprawl, especially critical in historically land-constrained Northeastern cities like New York and Philadelphia. However, identifying \acrshort{uvl}---and more broadly, developable land---is a manual, costly process that grows in complexity with local policy restrictions, environmental considerations, and density prioritization. An open-source tool to identify \acrshort{uvl} could facilitate redevelopment to housing, increase supply, and combat urban sprawl.

\begin{figure}[htbp]
  \centering
  \includegraphics[width=\textwidth]{building_permits.png}
  \caption{Building permits per capita since 2000. Data derived from the US Census Building Permits Survey. A sharp decrease in building permits is seen leading up to the Great Recession and permit volume has not yet recovered. Sun Belt MSAs such as Dallas--Fort Worth--Arlington, TX and Phoenix--Mesa--Chandler, AZ outpace land-constrained MSAs like New York--Newark--Jersey City, NY-NJ and Philadelphia--Camden--Wilmington, PA-NJ-DE-MD. Building permits, with the exception of NYC, still favor single-family one-unit dwellings (red), especially in the Sun Belt MSAs. New York--Newark--Jersey City is the only MSA where multi-family dwellings take precedence over single-family ones, but Philadelphia and other northeastern cities are catching up.}
  \label{fig:building_permits}
\end{figure}

\section{UVL to Mitigate Flood Risk}
\label{sec:flood_risk}

Urban expansion and the resulting increase in impervious surfaces create a cascade of urban water-related challenges: flood risk, sewer overflow, water quality deterioration and scarcity, and extreme weather~\cite{zhang2018Urbanization}. While coastal flooding risk from sea level rise is well modeled, the rapidly increasing threat of pluvial flooding---exacerbated by urban expansion and consequent soil sealing---is often overlooked and arguably more dangerous due to its low predictability. Existing flood risk projections tend to focus on coastal cities and frequently fail to account for pluvial flooding, which is significantly increasing risk in inland cities traditionally considered safe.

The heavy rains following Hurricane Ida in New York City exemplified this vulnerability (\Cref{fig:ida_heatmap}), where sewer systems were overwhelmed by intense, rapid-onset rain, resulting in widespread structural damage across inland urban areas---a notable deviation from the coastal areas typically affected during floods in NYC. Repurposing \acrshort{uvl} is not only a promising strategy for stormwater management but in many cases a necessary one. Left unmanaged, sealed \acrshort{uvl} can worsen water quality by allowing polluted runoff to enter surrounding systems. However, when redesigned using low-impact development strategies such as rain gardens and bioswales, \acrshort{uvl} can capture, slow, and filter stormwater at the source. Repurposing \acrshort{uvl} for stormwater facilities is a viable tool to enhance climate resilience in both coastal and inland cities, particularly against pluvial flooding. That said, effectiveness depends on site-specific conditions. For example, clay-heavy soils resist infiltration, increasing the effort required to convert \acrshort{uvl} to effective stormwater infrastructure. Ultimately, the effectiveness of \acrshort{uvl} repurposing depends on both design and location. Dispersed, low-connectivity vacant parcels in highly impervious areas can reduce peak runoff and support aging infrastructure, though local conditions such as soil composition may limit performance. Overall, \acrshort{uvl} offers a flexible, distributed approach to stormwater management and improving urban flood resilience.

\begin{figure}[htbp]
  \centering
  \includegraphics[width=0.85\textwidth]{DamagedBuildingsIdaHeatMap.pdf}
  \caption{Property damage from Cyclone Ida (2021), from the NYC Office of Management and Budget. While not as devastating as hurricanes Sandy and Irene, Cyclone Ida is particularly concerning as the damages were mainly from pluvial flooding due to rainfall, not coastal storm surge. Damages extended inland rather than concentrating in coastal areas---a record-breaking 3.15~inches/hour completely overwhelmed the aging sewer system designed for 1.75~in/hour.}
  \label{fig:ida_heatmap}
\end{figure}

\section{UVL for Social Good}
\label{sec:social_good}

While housing shortages and flood mitigation are two of the most urgent challenges that \acrshort{uvl} can help address, they are far from the only benefits of repurposing these spaces. \acrshort{uvl} already plays a significant social and ecological role in cities. In New York City, many vacant lots are not truly ``unused'' but instead function as community gardens, residential yards, parks, or recreational spaces~\cite{kremer2013social}. Beyond their existing uses, \acrshort{uvl} presents substantial opportunities for improving residents' wellbeing through urban greening. Converting abandoned spaces into green infrastructure---such as small urban farms or community shelters---has been shown to positively impact mental health, promote physical activity, and foster social interaction. Residents often perceive neglected vacant spaces as unsafe or unmanageable; however, their transformation into green, community-oriented areas can reduce stress, improve neighborhood perception, and strengthen social cohesion~\cite{jin2021Can}. Even in an unmanaged state, vacant lots can act as informal green spaces that support vegetation and urban wildlife, including avian species diversity~\cite{dennison2026role}. More intentional greening efforts can further enhance these ecological benefits while integrating \acrshort{uvl} into broader green infrastructure networks.

Overall, the identification and repurposing of \acrshort{uvl} offers a flexible, distributed approach to addressing both environmental and social challenges in cities. It can introduce elasticity into constrained housing markets, improve stormwater management, enhance biodiversity, and elevate quality of life. However, \acrshort{uvl} repurposing is not a one-size-fits-all solution; its success depends on careful, context-sensitive design that balances environmental goals with the needs and values of local communities.
